\chapter{Building Scholarly Structure}
\label{chap:structure}
\index{document structure}\index{abstract}\index{title metadata}\index{appendix}

Readers do not experience a scientific document as one uninterrupted stream.
They use a title to identify it, an abstract to judge relevance, headings to
navigate the argument, and appendices to find supporting detail. \LaTeX{} works
best when the source names these roles instead of imitating their appearance.

\begin{learningoutcomes}
  \item Create title metadata and print it with \cmd{maketitle}.
  \item Insert an abstract and keywords in a generic article.
  \item Organise content with sections and subsections.
  \item Generate a table of contents and explain repeat compilation.
  \item Add an appendix without typing its letter by hand.
  \item Distinguish logical structure from manual visual formatting.
\end{learningoutcomes}

\section{Metadata before appearance}

A title, author, and date describe the document. In the standard article class
we store them in the preamble and ask the class to typeset them:

\begin{latexcode}{Complete example: article title}
\documentclass{article}

\title{A Reproducible Comparison of Two Measurement Methods}
\author{Student Author}
\date{2026}

\begin{document}
\maketitle
This is a teaching document with simulated data.
\end{document}
\end{latexcode}

\sourceanatomy

\cmd{title}, \cmd{author}, and \cmd{date} store required arguments. They do not
print the title by themselves. \cmd{maketitle} in the body asks the class to
construct the title block from the stored information. An empty
\cmd{date}\required{} suppresses the date; \cmd{today} uses the compilation
date. For a versioned scholarly file, an explicit date is often easier to
audit.

\expectedoutput

The PDF begins with a centred title block using the article class's design,
followed by the teaching statement. The source does not specify font sizes or
vertical distances for each line.

\begin{tryit}
Change only the class from \code{article} to \code{report}. Compile and compare
the title treatment. Restore \code{article}. This demonstrates that metadata
can remain stable while the class changes broad design.
\end{tryit}

\section{The abstract}

An abstract is a compact account of purpose, approach, principal result, and
meaning. The standard \code{article} class provides an environment:

\begin{latexcode}{Abstract and keywords}
\begin{abstract}
We demonstrate a reproducible workflow for comparing two
measurement methods using simulated paired values. The example
is pedagogical and makes no empirical claim.
\end{abstract}

\noindent\textbf{Keywords:} agreement; measurement; reproducibility;
simulated data
\end{latexcode}

This fragment belongs after \cmd{maketitle} and before the first section. The
keyword line is an explicitly formatted convention for a generic article; the
standard class has no universal \cmd{keywords} command. Publisher classes often
provide their own command or environment. Use theirs when instructed.

\begin{goodpractice}
Do not put citations, undefined abbreviations, or conclusions unsupported by
the body into an abstract unless the target venue explicitly permits them.
Typesetting cannot repair an abstract that does not represent the study.
\end{goodpractice}

\section{Sections describe the argument}

Use the heading level that represents the relationship between ideas:

\begin{latexcode}{Complete example: section hierarchy}
\documentclass{article}
\begin{document}
\section{Introduction}
The study asks whether two methods agree.

\section{Methods}
\subsection{Study design}
Paired simulated values are used.

\subsection{Analysis}
Differences are calculated for each pair.

\section{Results}
Results will be added after the method is defined.
\end{document}
\end{latexcode}

\cmd{section} creates a major division; \cmd{subsection} creates a division
inside the current section; \cmd{subsubsection} is available when a third level
is genuinely needed. Do not skip from a section to a subsubsection simply to
obtain a smaller heading.

The starred form, such as \cmd{section*}, usually suppresses numbering and a
table-of-contents entry. It suits occasional unnumbered matter, not every
heading in a manuscript that requires navigation.

\begin{commonerror}
\textbf{Mistake:} making a heading with \cmd{textbf} and blank lines.

\textbf{Consequence:} it looks like a heading but has no structural level,
number, contents entry, or reliable target for a reference.

\textbf{Correction:} replace it with the appropriate section command and let
the class format it.
\end{commonerror}

\section{A table of contents}

The command \cmd{tableofcontents} prints headings recorded during compilation.
Place it where the contents should appear, commonly after front matter in a
long report or thesis.

On the first pass, \LaTeX{} writes heading information to a \file{.toc} auxiliary
file. On the next pass it can read that file. This is why a new section may not
appear in the contents until compilation repeats. \pkg{latexmk} notices the
change and runs the needed pass.

Short journal articles rarely include a table of contents. The command is
useful for reports, theses, books, and long internal documents. Structure and
reader need should decide, not the mere availability of a command.

\section{Appendices}

The command \cmd{appendix} changes the numbering of later top-level divisions.
In an article, later sections normally become Appendix A, Appendix B, and so
on.

\begin{latexcode}{Appendix pattern}
\appendix
\section{Additional checks}
This appendix records sensitivity checks for the simulated example.
\end{latexcode}

Do not type ``Appendix A'' inside the section title. \LaTeX{} supplies the letter
and updates it if another appendix is inserted earlier.

\begin{scienceinpractice}
An earth-science paper might place a detailed geological sample description in
an appendix while keeping the sampling rationale in Methods. The appendix is
still part of the scholarly record and should be cited from the main text,
labelled clearly, and reviewed with the same care.
\end{scienceinpractice}

\section{A complete structural skeleton}

The source below contains a document's architecture but very little prose. It
is useful as a planning skeleton because every heading represents a purpose.

\begin{latexcode}{Complete example: generic scholarly skeleton}
\documentclass[11pt]{article}

\title{A Reproducible Comparison of Two Measurement Methods}
\author{Student Author}
\date{2026}

\begin{document}
\maketitle

\begin{abstract}
This educational example uses simulated paired values.
\end{abstract}

\noindent\textbf{Keywords:} agreement; reproducibility;
simulated data

\section{Introduction}
State the problem and objective.

\section{Methods}
Describe data generation and analysis.

\section{Results}
Report the findings without interpretation.

\section{Discussion}
Interpret the demonstration and its limitations.

\section{Conclusion}
Answer the objective briefly.

\appendix
\section{Simulated data record}
Document how the teaching values were created.
\end{document}
\end{latexcode}

The imperative sentences are planning notes, not final manuscript prose. A
useful skeleton exposes missing logical components early; it should not tempt
us to fill headings with generic text that says nothing.

\begin{pausepredict}
Insert a new section before Results. Predict what happens to all later section
numbers after compilation. Now imagine those numbers had been typed into the
headings by hand.
\end{pausepredict}

\chapterproject

\begin{miniproject}
Turn the running article into the structural skeleton above. Replace
\code{Student Author} with \placeholder{Author name}; keep it visibly a
placeholder rather than inventing an identity. Retain the explicit simulated-
data statement. Compile twice or use \pkg{latexmk}. Confirm that the abstract,
keywords, five main sections, and appendix appear in the expected order.
\end{miniproject}

\chaptersummary

Title commands store metadata and \cmd{maketitle} prints it according to the
class. The abstract environment marks a scholarly role; generic keywords need
a documented local convention. Section commands express hierarchy and enable
automatic numbering and contents entries. \cmd{tableofcontents} may require a
repeat pass, and \cmd{appendix} changes later numbering without hard-coded
letters. Logical structure should lead visual formatting.

\chaptercheck
\begin{chapterchecklist}
  \item My title information is stored in the preamble.
  \item My abstract states what the demonstration does and does not claim.
  \item Every heading command represents a real level of the argument.
  \item I have not typed section or appendix numbers by hand.
  \item I understand why the contents may need another compiler pass.
  \item The running article now has a complete scholarly skeleton.
\end{chapterchecklist}

\chapterexercisestitle
\begin{chapterexercises}
  \item Explain the different jobs of \cmd{title} and \cmd{maketitle}.
  \item Add an explicit date to a minimal titled article.
  \item Convert a bold manual heading into semantic source.
  \item Predict the numbering after inserting a section before Methods.
  \item Explain why the generic article's keyword line is not universal.
  \item Add two automatically lettered appendices.
  \item Debug an abstract placed before \cmd{begin}\required{document}.
  \item Challenge: plan the heading hierarchy for a report in your field and
        justify every level.
\end{chapterexercises}
